Figure~\ref{fig:joint-weather-detail} shows that the upper-wind excess is not confined to one temperature range. Among the four temperature groups with adequate support, O/E at mean wind of at least 15 m/s ranges from 1.12 at 0--6$^{\circ}$C to 2.15 immediately below freezing. The corresponding high-to-low-wind contrasts range from 1.25 to 2.45, indicating that the magnitude of the wind pattern is not uniform across temperature groups; however, no formal interaction is estimated.

The elevation immediately below freezing is present in every wind category, while temperatures of at least 12$^{\circ}$C remain elevated in all three wind groups below 15 m/s. The cell combining mean wind of at least 15 m/s with temperature of at least 12$^{\circ}$C has the largest O/E (2.59), but it contains only 7 observed accidents compared with 2.70 expected and is the only sparse cell. It is therefore not interpreted as evidence of an especially strong combined effect. Overall, the joint results retain both the broad upper-wind excess and the non-linear temperature pattern.
